% Este arquivo é reutilizado em todos os modos para preservar a estrutura
% histórica do template. No modo projeto, ele apresenta Recursos e Cronograma;
% nos trabalhos acadêmicos finais, mantém o capítulo de Resultados.

\seprojeto{%
\chapter{RECURSOS E CRONOGRAMA}
\label{cap:recursos-cronograma}

\section{Recursos}
Descreva os recursos humanos, materiais, computacionais, laboratoriais ou financeiros necessários à execução da pesquisa. Adapte esta seção às exigências do curso, programa ou edital.

\section{Cronograma}
Apresente o cronograma de execução das atividades. Para cronogramas extensos, prefira o ambiente moderno \texttt{longabnttblr} do \texttt{tabularray-abnt}; \texttt{tabularx} e \texttt{longtable} continuam disponíveis para documentos legados. Evite reduzir o corpo da tabela abaixo do padrão do documento.

\begin{table}[htbp]
    \centering
    \UFCtab{
        \Caption{\label{tab:cronograma-exemplo} Exemplo simplificado de cronograma}
    }{
        \begin{abnttblr}[]{
            colspec={lcccc},
            row{1}={font=\bfseries},
            hline{1,Z}={1pt},
            hline{2}={0.6pt}
        }
            Atividade & Etapa 1 & Etapa 2 & Etapa 3 & Etapa 4 \\
            Revisão bibliográfica & X & X &  &  \\
            Desenvolvimento &  & X & X &  \\
            Experimentos &  &  & X & X \\
            Redação &  &  & X & X \\
        \end{abnttblr}
    }{
        \Fonte{Elaborada pelo autor.}
    }
\end{table}

}{%
\chapter{Resultados}
\label{chap:resultados}

Texto texto texto texto texto texto texto texto texto texto texto texto texto texto texto texto texto texto texto texto texto texto texto texto texto texto texto texto texto texto texto texto texto texto texto texto texto texto texto texto texto texto texto texto texto texto texto texto texto texto texto texto texto texto texto texto texto texto texto texto texto texto texto texto texto texto texto texto texto.

\section{Resultados do Experimento A}
\label{sec:resultados-do-experimento-a}

Procure deixar as figuras dos resultados o maior possível preenchendo a largura do texto do documento que possui $16~cm$.

\begin{figure}[htbp]
    \centering
    \UFCfig{
        \Caption{\label{fig:tensaoimpedanciahumana} Gráfico de tensão considerando a impedância humana}
    }{
        \fbox{\UFCincludegraphics[width=.96\linewidth]{tensaoimpedanciahumana}}
    }{
        \Fonte{Elaborada pelo autor.}
    }
\end{figure}

Texto texto texto texto texto texto texto texto texto texto texto texto texto texto texto texto texto texto texto texto texto texto texto texto texto texto texto texto texto texto texto texto texto texto texto texto texto texto texto texto texto texto texto texto texto texto texto texto texto texto texto texto texto texto texto texto texto texto texto texto texto texto texto texto texto texto texto texto texto.

\begin{figure}[htbp]
    \centering
    \UFCfig{
        \Caption{\label{fig-grafico-1} Produção anual das dissertações de mestrado e teses de doutorado entre 1990 e 2008}
    }{
        \fbox{\UFCincludegraphics[width=.96\linewidth]{figura-3}}
    }{
        \Fonte{Elaborada pelo autor.}
    }
\end{figure}

Texto texto texto texto texto texto texto texto texto texto texto texto texto texto texto texto texto texto texto texto texto.

Texto texto texto texto texto texto texto texto texto texto texto texto texto texto texto texto texto texto texto texto texto texto texto texto texto texto texto texto texto texto texto texto texto texto texto texto texto texto texto texto texto texto texto texto texto texto texto texto texto texto texto texto texto texto texto texto texto texto texto texto texto texto texto texto texto texto texto texto texto.

\section{Resultados do Experimento B}
\label{sec:resultados-do-experimento-b}

Texto texto texto texto texto texto texto texto texto texto texto texto texto texto texto texto texto texto texto texto texto texto texto texto texto texto texto texto texto texto texto texto ..

\begin{table}[htbp]
    \centering
    \UFCtab{
        \Caption{\label{tab:notas} Notas dos participantes nas avaliações A, B e C}
    }{
        \begin{tabular}{l S[table-format=2.0] S[table-format=2.0] S[table-format=2.0]}
            \toprule
            Identificação & {Avaliação A} & {Avaliação B} & {Avaliação C} \\
            \midrule
            Participante 1 & 7 & 9 & 10 \\
            Participante 2 & 8 & 2 & 1 \\
            Participante 3 & 5 & 10 & 6 \\
            Participante 4 & 3 & 1 & 4 \\
            Participante 5 & 2 & 4 & 1 \\
            Participante 6 & 0 & 7 & 2 \\
            \bottomrule
        \end{tabular}
    }{
        \Fonte{Elaborada pelo autor.}
    }
\end{table}

 Texto texto Referenciando a \autoref{tab:notas}  texto texto texto texto texto texto texto texto texto texto texto texto texto texto texto texto texto texto texto texto texto texto texto texto texto texto texto texto texto texto.Texto texto texto texto texto texto texto texto texto texto texto texto texto texto texto texto texto texto texto texto texto.

Texto texto texto texto texto texto texto texto texto texto texto texto texto texto texto texto texto texto texto texto texto texto texto texto texto texto texto texto texto texto texto texto texto texto texto texto texto texto texto texto texto texto texto texto texto texto texto texto texto texto texto texto texto texto texto texto texto texto texto texto texto texto texto texto texto texto texto texto texto.Texto texto texto texto texto texto texto texto texto texto texto texto texto texto texto texto texto texto texto texto texto texto texto texto texto texto texto texto texto texto texto texto texto texto texto texto texto texto texto texto texto.

Texto texto texto texto texto texto texto texto texto texto texto texto texto texto texto texto texto texto texto texto texto texto texto texto texto texto texto texto texto texto texto texto texto texto texto texto texto texto texto texto texto texto texto texto texto texto texto texto.Texto texto texto texto texto texto texto texto texto texto texto texto texto texto texto texto texto texto texto texto texto.

Texto texto  Referenciando a \autoref{tab:notas}  texto texto texto texto texto texto texto texto texto texto texto texto texto texto texto texto texto texto texto texto texto texto texto texto texto texto texto texto texto texto texto texto texto texto texto texto texto texto texto texto texto texto texto texto texto texto texto texto texto texto texto texto texto texto texto texto texto texto texto texto texto texto texto texto texto texto texto.
}
